\exercisesection{4}

\begin{enumerate}
\def\labelenumi{\arabic{enumi}.}
\tightlist
\item
  \begin{enumerate}
  \def\labelenumii{(\alph{enumii})}
  \tightlist
  \item
    设 G 为全序群，M 为 \(G_{+}\) 中不含 0 的优势集。证明：存在 G 中与 M 不相交的最大孤立子群 H。
  \end{enumerate}
\end{enumerate}

\begin{enumerate}
\def\labelenumi{(\alph{enumi})}
\setcounter{enumi}{1}
\item
  设 A 为赋值环，v 为与 A 相伴的赋值。A 的理想 \(p \neq 0\) 为素理想的充要条件是：p 对应于 v 的全序取值群 G 中的一个优势集 M，并且若 H 为 G 中与 M 不相交的最大孤立子群，则 M 是 \(H_{+}\) 在 \(G_{+}\) 中的补集。
\item
  沿用 (b) 的记号，A 的理想 q 为 p-准素理想的充要条件是满足下列条件之一：或者 \(H = \{0\}\)（此时 p 为极大理想），或者 q = p。
\item
  对于 \(\mathfrak{a}\) 这一 \(\mathbf{A}\) 的理想，若它既不等于 0 也不等于 \(\mathbf{A}\)，则满足 \(\mathfrak{a}^2 = \mathfrak{a}\) 的充要条件是：\(\mathfrak{a}\) 对应于一个优势集 \(\mathbf{M}\)，且若 \(\mathbf{H}\) 是 \(\mathbf{G}\) 中与 \(\mathbf{M}\) 不相交的最大孤立子群，则 \(\mathbf{M}'\) 是 \(\mathbf{M}\) 在全序群 \(\mathbf{G}' = \mathbf{G} / \mathbf{H}\) 中的像，且没有最小元；此时 \(\mathfrak{a}\) 是素理想。
\end{enumerate}

\begin{enumerate}
\def\labelenumi{\arabic{enumi}.}
\setcounter{enumi}{1}
\tightlist
\item
  设 G 为全序群。
\end{enumerate}

\begin{enumerate}
\def\labelenumi{(\alph{enumi})}
\item
  对于 G 的每个元素 a \textgreater{} 0，令 \(\mathrm{H}(a)\) 为包含 a 的最小孤立子群，令 \(\mathrm{H}^{-}(a)\) 为不包含 a 的最大孤立子群；证明 \(\mathrm{H}(a)/\mathrm{H}^{-}(a)\) 同构于 R 的一个子群。
\item
  反之，若 H 是 G 的孤立子群，且包含于 H 的 G 的孤立子群所成集合中存在关于包含关系的最大元 \(H'\) \(\neq H\)，则 \(\mathrm{H} = \mathrm{H}(a)\) 且 \(\mathrm{H}' = \mathrm{H}^{-}(a)\)，对所有 \(a \in H_{+} \cap CH_{+}'\) 均成立。这样的孤立子群 H 称为主孤立子群，\(H'\) 称为它的前驱。
\item
  若 \(\mathbf{H}_1, \mathbf{H}_2\) 是 \(\mathbf{G}\) 的两个不同孤立子群，且 \(\mathbf{H}_1 \subset \mathbf{H}_2\)，证明存在 \(\mathbf{H}, \mathbf{H}'\) 的两个孤立子群 \(\mathbf{G}\)，使得 \(\mathbf{H}_1 \subset \mathbf{H}' \subset \mathbf{H} \subset \mathbf{H}_2\)，并且 \(\mathbf{H} / \mathbf{H}'\) 同构于一个非零的 \(\mathbf{R}\) 子群。
\item
  令 \(\Sigma\) 为 \(G\) 的孤立子群所成的集合（按包含关系全序），令 \(\Theta \subset \Sigma\) 为主孤立子群所成的集合。证明 \(\Sigma\) 同构于 \(\Theta\) 的完备化（集合论，第三章 § 1，习题 15）。
\end{enumerate}

\begin{enumerate}
\def\labelenumi{\arabic{enumi}.}
\setcounter{enumi}{2}
\tightlist
\item
  \begin{enumerate}
  \def\labelenumii{(\alph{enumii})}
  \tightlist
  \item
    令 \(\Theta\) 为全序集，并对每个 \(s \in \Theta\) 令 \(\mathbf{A}_s\) 为不退化为 0 的 \(\mathbf{R}\) 的子群。令
  \end{enumerate}
\end{enumerate}

\[
\Gamma (\Theta , (A _ {s}) _ {s \in \Theta}) = G
\]

是由支集为 \(\biguplus_{s\in\Phi}\mathbf{A}_{s}\) 的良序子集的函数组成的乘积 \(f\) 的子群。我们在 \(\Theta\) 上定义一个与其群结构相容的序：令元素 \(\mathbf{G}\) 所成的集合 \(\mathbf{G}_{+}\) 由函数 0 以及满足下述条件的函数 \(\geqslant0\) 组成：f 在其支集的最小元 \(f\neq0\) 处满足 \(f(\theta)>0\)。证明 \(\theta\) 是全序群，并且 \(f\) 按关系 \(\mathbf{G}\) 排序后典范同构于 \(\Theta\) 的主孤立子群所成的集合；此外，在此同构下若 \(\mathbf{G}\) 对应于 \(\supset\)，则 \(\mathrm{H}(a)\) 同构于 \(s\in\Theta\)。\(\mathrm{H}(a)/\mathrm{H}^{-}(a)\)\(\mathbf{A}_{s}\)

\begin{enumerate}
\def\labelenumi{(\alph{enumi})}
\setcounter{enumi}{1}
\tightlist
\item
  考虑全序群 \(G'\)（按字典序排序）的子群 \(Q \times Q\)，它由元素 \((p_{n}^{-1}, np_{n}^{-1})\) 生成，其中 \((p_{n})\) 是严格递增的素数序列。证明 G' 中既不同于 0 也不同于 \(H'\) 的唯一孤立子群 \(G'\) 是 \(G'\)；但 G' 不同构于按字典序排序的乘积 \(\{0\} \times Z\)\(G'\)\(H' \times (G'/H')\) 的任何子群。
\end{enumerate}

\begin{enumerate}
\def\labelenumi{\arabic{enumi}.}
\setcounter{enumi}{3}
\tightlist
\item
  设 \(A\) 为非域的赋值环。
\end{enumerate}

\begin{enumerate}
\def\labelenumi{(\alph{enumi})}
\item
  证明 A 为离散赋值环的充要条件是 A 的每个素理想都是主理想。
\item
  证明：若 A 是高度 1 的赋值的环，则 A 的分式域是有限生成的 A-代数。
\end{enumerate}

* 5. 设 K 为域，M 为 K 中非域子环所成的集合。证明 M 的极大元是其分式域 L 的一个高度 1 的赋值环，并且 K 是 L 的代数扩张。下列条件等价：

(α) M 非空。

(β) M 有极大元。

(γ) K 上存在高度 1 的赋值

(δ) K 不是有限素域的代数扩张（参见 § 8，第 1 节）。*

¶ 6. (a) 设 S 为没有孤立点的超 Stonean 紧空间（积分论，第五章 § 5，习题 14）。令 \(\mathcal{C}_{0}(S)\) 为 S 上取实值（有限或非有限）、在 S 上连续且满足集合 \(\overline{f}(-\infty)\) 与 \(\overline{f}(+\infty)\) 在 S 中处处不稠密的函数所成的向量空间（积分论，第二章 §1，习题 13 (g)）。证明：对 S 上每个测度 \(\mu > 0\)，\(\mathcal{C}_{0}(S)\) 中存在上积分为无穷的函数 f（注意，对于 S 的每个处处不稠密闭子集 N，都存在 \(\mathcal{C}_{0}(S)\) 中在 N 上等于 \(+\infty\) 的函数 f；然后证明只需考虑测度 \(\mu\) 为正规测度的情形，并将 f 适当地定义为 \(\mathcal{C}_{0}(S)\) 中一列函数的上界）。

\begin{enumerate}
\def\labelenumi{(\alph{enumi})}
\setcounter{enumi}{1}
\item
  令 \(\mathbf{G}\) 为 \(\mathcal{C}_0(S)\) 的有序子群，它由 \(\mathcal{C}_0(S)\) 中取值于 \(\mathbf{Z}\) 或 \(\pm \infty\) 的函数组成。证明 \(\mathbf{G}\) 是完备格，并且（作为有序群）不同构于乘积群 \(\mathbf{R}^{\mathbf{l}}\) 的任何子群。
\item
  令 K 为域， 为以空间 S 为指标集的一族不定元（\(A_{0} = K[[X_{s}]]_{s \in S}\)）的系数在 K 中的形式幂级数环（代数，第 IV 章 § 5，习题 1）。令 b 为满足下述性质的形式幂级数 \(X_{s}\) 所成的集合：存在 S 的处处不稠密子集 N，使得 P 的每个系数 \(P \in A_{0}\) 的单项式都含有某个满足 \(N \subset S\)\(\neq 0\) 的 。证明 b 是 \(X_{s}\) 的理想，并且环 \(s \in N\)\(A_{0}\) 是整环且完全整闭（证明后一点时应用第五章 § 1，第 4 节的命题 14，并仿照第五章 § 1，习题 10 论证；还要使用 S 中每个第一类集都是处处不稠密的事实）。\(A = A_{0}/b\) 是整环且完全整闭（证明后一点时应用第五章 § 1，第 4 节的命题 14，并仿照第五章 § 1，习题 10 论证；还要使用 S 中每个第一类集都是处处不稠密的事实）。
\item
  令  为群 \(f\) 中满足  的元素，令 N 为 \(\geqslant 0\) 的点所成的处处不稠密集合。令 \(G\)\(N\)\(f(s) = \pm \infty\) 为形式幂级数 \(p_f\) 在 A 中的像；映射 \(A\)\(\prod_{s \notin N} (1 + X_s)^{f(s)}\) 是从加法幺半群 \(f \mapsto p_f\) 到 A 的一个乘法子幺半群的同构。推出 A 不可能是其分式域的高度 1 的一族赋值环的交（证明这将蕴含 G 同构于乘积 \(G_+\) 的一个子群）。（参见习题 8、9。）\(A\)\(A\)\(G\)\(R^1\) 的一个子群）。（参见习题 8、9。）
\end{enumerate}

¶ 7. 若局部整环 A 不是域，且 A 除 (0) 与极大理想 m 外没有别的素理想，则称 A 的维数为 1。对 A 的每个理想 a，令 A: a 表示 A 的分式域中 a 在 A 中的传递子所成的 A-子模（它总是包含 A）。以下均假设 A 的维数为 1。

\begin{enumerate}
\def\labelenumi{(\alph{enumi})}
\item
  证明：若 A 是强 Lasker 环（第四章 § 2，习题 28），则 A: m ≠ A。（反设相反，注意对每个整数 r \(^{r}\) 0 都有 A: m \textgreater{} = A；另一方面，A 中包含于 m 的每个理想 ≠ 0 都是 m-准素的，因而特别地，每个主理想 Ax ⊆ m（其中 x ≠ 0）都是 m-准素的；最后注意 h ≠ k 时 m \(^{h}\) ≠ m \(^{k}\)（第四章 § 2，习题 29 (d)）。）
\item
  证明：对于维数为 1 的强 Lasker 局部整环 A，下列性质等价：
\end{enumerate}

\((\alpha)\) A 是完全整闭的。

(β) 极大理想 m 是主理想。

(γ) 每个 m-准素理想都形如 \(m^{k}\)。

(δ) 理想 m 是可逆的（第二章 § 5，第 6 节），这等价于 m. (A: m) = A。

(ε) A 是离散赋值环。

（为证明  蕴含 \((\alpha)\)，使用 (a)，并注意若对所有 \((\delta)\) 都有 ，则 \(\mathfrak{m}.\left(\mathbf{A}:\mathfrak{m}\right)^k = \mathfrak{m}\) 不是完全整闭的（参见第七章 §1，习题 4）。为证明 \(k > 0\)\(\mathbf{A}\) 蕴含 \((\delta)\)，注意对于每个包含于 \((\gamma)\) 的理想 ，都可写成 \(\mathfrak{a} \neq 0\)，其中 \(\mathfrak{m}\)\(\mathfrak{a} = \mathfrak{ma}_1\)，然后使用第四章 §2，习题 29 (d)。最后，为证明 \(\mathfrak{a}_1 \neq 0\) 蕴含 \((\gamma)\)，注意 \((\beta)\)，且 \(\mathfrak{m} \neq \mathfrak{m}^2\) 蕴含 \((\gamma)\) 是一个 1 维的 \(\mathfrak{m} / \mathfrak{m}^2\)-向量空间。）\((\mathbf{A} / \mathfrak{m})\)-向量空间。）

\begin{enumerate}
\def\labelenumi{(\alph{enumi})}
\setcounter{enumi}{2}
\item
  证明：用第三章 § 3，习题 15 (b) 的记号，局部环 \(\mathbf{A}_{\mathfrak{m}}\) 是维数为 1 的 Noether 整环，但不是整闭的。
\item
  设 K 为代数闭域，A 为域 \(\mathbf{K}(\mathbf{X},\mathbf{Y})\) 的子环；该域是 K 上两个不定元的有理函数域，而 A 由满足下述条件的元素 \(f\in\mathbf{K}(\mathbf{X},\mathbf{Y})\) 组成：可将 X 代入 f 中的 0，且 \(f(0,\mathbf{Y})\) 属于 K。证明 A 是强 Lasker 的整闭局部整环，维数为 1，但不是完全整闭的。（若 B 是多项式环 \(\mathbf{K}[\mathbf{X},\mathbf{Y}]\)，p 是 B 的素理想 BX，注意 A 是 \(B_{p}\) 中等于 \(K+pB_{p}\) 的子环。）(*)
\end{enumerate}

¶ 8. 设 A 为维数 1 的局部整环（习题 7），K 为其分式域。证明 A 是（K 的）高度 1 的赋值环的交的充要条件是 A 完全整闭。依次证明：

\begin{enumerate}
\def\labelenumi{(\alph{enumi})}
\item
  含有 A 以及 A 的极大理想 m 中一个元素 \(\neq0\) 的逆元的 K 的子环等于 K（使用 A 中每个既不等于 0 也不等于 A 的理想都是 m-准素理想这一事实）。
\item
  若 \(\mathbf{B} \supset \mathbf{A}\) 是不同于 \(\mathbf{K}\) 的 \(\mathbf{K}\) 的子环，则存在高度 1 的赋值环 \(\mathbf{V}\)，使得 \(\mathbf{B} \subset \mathbf{V}\)（使用 (a)）。
\item
  令 \(z \in \mathbf{K} - \mathbf{A}\)；仿照习题 7 (b) 论证 \(zm \subset m\) 不可能。推出存在高度 1 的赋值环 \(V\)，使得 \(A \subset V\) 且 \(z \notin V\)。（若 \(zm \subset A\)，使用 (a)；否则存在 \(x \in m\)，使得
\end{enumerate}

\[
x y = z \in \mathrm{K} - \mathrm{A};
\]

注意 \(z \notin A[z^{-1}]\)，并使用 (b) 证明 V 的存在性。）

¶ 9. 设 A 为 Lasker 完全整闭整环（第四章 §2，习题 23）。进一步假设，对 A 中所有 ，与 A/Ax 弱伴随的素理想 \(x \neq 0\)（第四章 § 1，习题 17）全都具有高度 1，也就是说对所有 i，\(\mathfrak{p}_i\) 不含除自身与 0 以外的素理想（参见第七章 § 1，第 6 节）。证明在这些条件下 A 是高度 1 的赋值环的交。（先使用第四章 § 1 的习题 17 (i)，证明 A 是所有环  的交，其中 p 遍历高度 1 的素理想所成集合。然后使用第四章 § 2，习题 23 的条件（），证明这些局部环  完全整闭。最后应用习题 8。）当 A 是每个理想都具有单个准素分解的 Lasker 完全整闭整环时，上述假设特别得到满足（第四章 § 2，习题 26）。\(p_{i}\)\(A_{p}\)\(A_{p}\)\(LA_{I}\) 完全整闭。最后应用习题 8。）当 A 是每个理想都具有单个准素分解的 Lasker 完全整闭整环时，上述假设特别得到满足（第四章 § 2，习题 26）。

¶ 10. 设 A 为一个主理想集合按包含关系全序的环。

\begin{enumerate}
\def\labelenumi{(\alph{enumi})}
\item
  证明 \(\mathbf{A}\) 是局部环，且其中的幂零根 \(\mathfrak{N}\) 是素理想。证明或者 \(\mathfrak{N}^2 = \mathfrak{N}\)，或者 \(\mathfrak{N}\) 是幂零的；环 \(\mathbf{A} / \mathfrak{N}\) 是赋值环。
\item
  若 \(\mathcal{L}\) 是 A 的主理想所成的全序集，证明 A 的理想所成的集合按包含关系全序，并同构于集合 \(\mathcal{L}\) 的完备化（集合论，第三章 § 1，习题 15）。
\item
  假设 ；于是可将 \(\mathfrak{N}^2 = 0\) 看作赋值环 \(\mathfrak{N}\) 上的模；令 K 为 V 的分式域，令 \(V = A / \mathfrak{N}\) 为 K 上与 V 对应的赋值 v 的阶群；证明 \(K\)\(V\)\(\Gamma\) 中存在两个优势集 \(v\)\(K\)\(V\)\(\Gamma\)，使得 \(M \subset M'\) 作为 V-模同构于 \(\mathfrak{N}\)（§3，第 4 节的记号）。\(V\)\(\mathfrak{a}(M') / \mathfrak{a}(M)\)（§3，第 4 节的记号）。
\item
  反之，给定分式域为 K、阶群为  的赋值环 V，以及 \(\Gamma\) 中的两个优势集 ，令 \(M \subset M'\)，并在乘法群 \(\Gamma\)\(Q = a(M') / a(M)\) 上规定乘法为\(A = V \times Q\) 上规定乘法为
\end{enumerate}

\[
(z, t) (z ^ {\prime}, t ^ {\prime}) = (z z ^ {\prime}, z t ^ {\prime} + z ^ {\prime} t);
\]

证明在 A 中主理想所成集合按包含关系全序；有序对 \((0, t)\)（其中 \(t \in Q\)）所成的集合 N 是 A 的幂零根，\(N^{2} = 0\)，且 A/N 同构于 V。

\begin{enumerate}
\def\labelenumi{(\alph{enumi})}
\setcounter{enumi}{4}
\item
  对每个素数 ，环 \(p\) 的主理想所成集合按包含关系全序，且 A 的幂零根 \(A = \mathbf{Z} / p^2\mathbf{Z}\) 满足 \(A\)\(\mathfrak{N}\)；但证明环 A 不同构于由 (d) 的方法（从 \(A\)\(\mathfrak{N}^2 = 0\) 与 \(A\)\(V = A / \mathfrak{N}\) 出发）构造的环（注意 A 不含同构于 \(Q = \mathfrak{N}\) 的子环）。\(A\)\(Z / p\mathbf{Z}\) 的子环）。
\item
  证明对于每个素数 ，群 \(p\) 关于素域 \(\mathbf{A}\)\(\mathbf{U}_p\) 的代数 \(\mathbf{F}_p\)（代数，第 VII 章 §2，习题 3）是一个环，其中主理想所成集合按包含关系全序，且幂零根  满足 \(\mathfrak{N}\)，但 \(\mathfrak{N}^2 = \mathfrak{N}\) 不同构于赋值环关于某个理想的商环。\(\mathbf{A}\) 不同构于赋值环关于某个理想的商环。
\end{enumerate}

\begin{enumerate}
\def\labelenumi{\arabic{enumi}.}
\setcounter{enumi}{10}
\tightlist
\item
  设 K 为带有高度 1 赋值 v 的域。
\end{enumerate}

\begin{enumerate}
\def\labelenumi{(\alph{enumi})}
\tightlist
\item
  令  为 \(\mathbf{P}(\mathbf{X}) = a_0\mathbf{X}^n +a_1\mathbf{X}^{n - 1} + \dots +a_n\) 中次数为 \(\mathrm{K[X]}\) 且满足 \(n > 1\) 的多项式；证明存在区间 \(a_{n} \neq 0\) 中严格递增的整数序列 ，使得：(1) \((i_{k})_{0 \leqslant k \leqslant r}\)，\([0, n]\)\(i_{0} = 0\)；(2) 对 \(i_{r} = n\)， 有限；(3) 对每个不同于各 \(v(a_{i_{k}})\) 的指标 ，若 \(0 \leqslant k \leqslant r\)\(j\) 且 \(0 \leqslant j \leqslant n\) 有限，则该点\(i_{k}\)\(v(a_{j})\) 有限，则该点
\end{enumerate}

\[
(j, v (a _ {j})) \in \mathbf {R} ^ {2}
\]

位于经过点 \((i_{k}, v(a_{i_{k}}))\) 与 \((i_{k+1}, v(a_{i_{k+1}}))\) 的直线之上；若 \(j < i_{k}\) 或 \(j > i_{k+1}\)，则严格位于该直线之上。连接点 \((i_{k}, v(a_{i_{k}}))\) 与 \((i_{k+1}, v(a_{i_{k+1}}))\) 的线段的并称为 P 的 Newton 多边形，上述线段称为边，点 \((i_{k}, v(a_{i_{k}}))\) 称为多边形的顶点。

\begin{enumerate}
\def\labelenumi{(\alph{enumi})}
\setcounter{enumi}{1}
\item
  假设 P 的所有零点都属于 K。证明 P 的所有零点具有相同赋值的充要条件是 r = 1（换言之，Newton 多边形退化为单条边）。（为证明条件充分，考虑乘积 \(P_{1}P_{2}\) 的 Newton 多边形，其中 \(P_{1}\) 的所有零点在 v 的环中均可逆，而 \(P_{2}\) 的所有零点都属于 v 的理想。）
\item
  假设 \(\mathbf{P}\) 的所有零点都属于 \(\mathbf{K}\)；作出 \(\mathbf{P}\) 的 Newton 多边形并令
\end{enumerate}

\[
\rho_ {k} = i _ {k + 1} - i _ {k}, \quad \sigma_ {k} = (v (a _ {i _ {k + 1}}) - v (a _ {i _ {k}})) / \rho_ {k}.
\]

证明：对 ，P 恰有 \(0 \leqslant k \leqslant r - 1\) 个零点（按重数计），其赋值都等于 \(\rho_k\)（使用 (b)，并对 r 作归纳）。\(\sigma_k\)\(r\)（使用 (b)，并对 r 作归纳）。

\begin{enumerate}
\def\labelenumi{(\alph{enumi})}
\setcounter{enumi}{3}
\tightlist
\item
  推广到任意赋值 v 的情形（将 v 的阶群 \(\Gamma\) 嵌入向量 Q-空间 \(\Gamma_{(Q)}\) 中，后者具有自然的全序群结构）。
\end{enumerate}

